% econ-slides speaker script for talk.tex (Presentation 2, progress report).
% One \scriptframe per main-deck frame, keyed to the exact frame title.
% Check: python3 <skill-dir>/scripts/check_script.py script.tex --deck talk.tex --slot-minutes 10
\documentclass[11pt]{article}
\usepackage[margin=0.78in,headheight=14pt]{geometry}
\usepackage{xcolor}
\usepackage{microtype}
\usepackage{needspace}
\usepackage{setspace}
\usepackage{fancyhdr}
% Shared title metadata and result phrases for talk.tex and script.tex.
% Every magnitude below is sourced from progress_brief.tex (P2's results
% write-up); background/model phrasing is sourced from research_brief.tex.
% Numbers are never rounded further for the slide than they are in the brief.

\newcommand{\PaperTitle}{Network Position and Academic Performance\\[0.15em]
  {\normalsize Adding Network Centrality to Smirnov \& Thurner (2017)}}
\newcommand{\PaperShortTitle}{Network Position and Academic Performance}
\newcommand{\PaperAuthors}{Amy Bing \and Tingying Huang}
\newcommand{\AuthorShort}{Bing \& Huang}
\newcommand{\PresenterName}{Tingying Huang}
\newcommand{\CoauthorsSpoken}{Amy Bing}
\newcommand{\PaperAffiliations}{Monash University}
% Shared affiliation for both authors -- the simple default path (no
% per-author \AuthorAffil tabular needed since both are Monash).
\newcommand{\TitleAuthorBlock}{\PaperAuthors}
\newcommand{\TitleAffiliationBlock}{\PaperAffiliations}
\newcommand{\KeyQuestion}{Does network position predict grades?}
\newcommand{\KeyQuestionSubtitle}{On top of who a student's direct friends are}
\newcommand{\VenueDate}{ECC3841 Network Economics}
\newcommand{\VenueShort}{Presentation 2 --- Progress Report}
\newcommand{\TitleDisclaimer}{}

% ---- Result phrases: slide (written) and spoken (script) forms share the ----
% ---- same source-traced number, worded for the medium.                  ----

% Result 1 -- replication. Source: progress_brief.tex, "Result 1".
\newcommand{\ResultOneSlide}{Own past GPA strongly predicts next GPA ($\beta=0.955$); a friend's past GPA does not (\(p=0.557\))}
\newcommand{\ResultOneSpoken}{a student's own past GPA is a very strong predictor of their next GPA, but once you control for that, a friend's past GPA predicts nothing}

% Result 2 -- naive effect is popularity. Source: progress_brief.tex, "Result 2".
\newcommand{\ResultTwoSlide}{Katz-Bonacich alone: $+0.093$ ($p<0.001$). Add degree: $-0.028$ ($p=0.446$)}
\newcommand{\ResultTwoSpoken}{on its own, centrality looks like it matters, but the moment you add raw popularity to the same regression, that effect disappears}

% Result 3 -- pooled headline, read at face value. Source: progress_brief.tex, "Result 3".
\newcommand{\ResultThreeSlide}{$\beta=-0.0415$, SE\(=0.019\), \(p=0.033\), \(n=36{,}696\)}
\newcommand{\ResultThreeSpoken}{a small, negative, statistically significant effect, at the value of phi we fixed in advance}

% Result 4 -- corrected test. Source: progress_brief.tex, "Result 4".
\newcommand{\ResultFourSlide}{No \(\phi\) is significant ($p=0.36$ to $0.93$); the sign is not even consistent}
\newcommand{\ResultFourSpoken}{at every value of phi, not significant, and the sign flips back and forth}

% Headline / conclusion. Source: progress_brief.tex, boxed conclusion.
\newcommand{\ResultHeadlineSlide}{Net of popularity, friend GPA, and a student's own past performance, we find no evidence that network position independently predicts grades}
\newcommand{\ResultHeadlineSpoken}{once we properly control for popularity, friend GPA, and a student's own past performance, there is no evidence that network position predicts grades on its own}
\newcommand{\ResultSupportSlide}{Naive centrality looked positive, the pooled sweep looked negative, the corrected test looks like nothing}
\newcommand{\ResultSupportSpoken}{the naive test looked positive, the pooled test looked negative, and the one correctly specified test looks like nothing at all}
\newcommand{\ResultImplicationSlide}{Interventions built around widening students' network reach are not supported by this data once the obvious confounds are handled}
\newcommand{\ResultImplicationSpoken}{so policies that try to help students by widening their network reach are not supported by what we've found here}


\definecolor{ScriptBlue}{HTML}{1F4E79}
\definecolor{ScriptGray}{HTML}{666666}
\setstretch{1.12}
\setlength{\parindent}{0pt}
\setlength{\parskip}{0.45em}
\pagestyle{fancy}
\fancyhf{}
\lhead{\small\textcolor{ScriptGray}{\PaperShortTitle}}
\rhead{\small\textcolor{ScriptGray}{Rehearsal script}}
\cfoot{\thepage}

\newcommand{\Click}[1]{\textcolor{ScriptBlue}{\textbf{[Click #1]}}}
\newcommand{\Pause}{\textcolor{ScriptGray}{\textit{[pause]}}}
\newcommand{\Cue}[1]{\par\nobreak\smallskip\noindent\textcolor{ScriptGray}{\footnotesize
  \textbf{Live cue:} #1}\par}

\newcommand{\scriptopening}[2]{%
  \Needspace{10\baselineskip}%
  \section*{Opening \hfill\normalfont\small\textcolor{ScriptGray}{#1}}
  #2\par}
\newcommand{\scriptframe}[3]{%
  \Needspace{9\baselineskip}%
  \par\medskip\noindent\textcolor{ScriptBlue}{\rule{\linewidth}{0.55pt}}%
  \par\smallskip\noindent
  \textbf{#1}\hfill\textcolor{ScriptGray}{\small #2}\par\smallskip
  #3\par}
\newcommand{\ScriptAppendix}{%
  \clearpage
  \section*{\mbox{Q\&A appendix} \hfill
    \normalfont\small\textcolor{ScriptGray}{conditional; excluded from talk time}}}

\title{Speaker script\\[0.25em]\large \PaperShortTitle}
\author{\PaperAuthors}
\date{\VenueDate\ --- \VenueShort\quad 10 minutes}

\begin{document}
\maketitle
\thispagestyle{fancy}

\scriptopening{about 0.4 min}{%
  Good morning. I'm Tingying Huang, presenting joint work with Amy Bing on
  network position and academic performance. Last time we asked whether a
  student's position in their friendship network predicts their grades, on
  top of who their direct friends are. Today is about what actually happened
  when we ran that test --- including a gap in our own design that we found
  and had to fix.
  \Cue{Greeting; coauthor; recap the question in one sentence; frame today as results, not setup.}
}

\scriptframe{\KeyQuestion}{about 0.9 min}{%
  Quick recap of where the question comes from. Smirnov and Thurner showed
  that friends' grades move together mostly because students pick friends
  with similar GPA. It's not that friends pull each other up or down. But
  their method only looks at one number, a friend's average GPA, and it
  ignores the rest of the network. Calv{\'o}-Armengol, Patacchini and Zenou
  showed that network position itself, measured with Katz-Bonacich
  centrality, predicts achievement. But their data is a single snapshot. So
  they can't tell whether position raises grades, or good grades just
  attract more friends. Our plan combines the two: take the centrality
  prediction to Smirnov and Thurner's panel data, and test it while holding
  friend-group GPA and a student's own past GPA fixed.
  \Cue{Two antecedents in one sentence each; the gap each leaves; our plan combines them.}
}

\scriptframe{Recap: design and data}{about 0.8 min}{%
  Here's the regression we planned. Next term's GPA depends on Katz-Bonacich
  centrality, on the student's own past GPA, on their friends' average GPA,
  and on raw in- and out-degree, plus group fixed effects. Friend GPA nets
  out friend-group composition. Degree nets out raw popularity. And we sweep
  phi from near zero, where centrality is basically popularity, to near its
  ceiling, where it captures long reach through the network. One fact about
  the data matters a lot for what comes later. GPA is measured six times over
  for the school group. But for each of the four university cohorts, it's a
  single static number, not something that varies over time.
  \Cue{Regression equation; each control's job; phi dial; the panel-vs-static data fact, stated plainly.}
}

\scriptframe{What we did}{about 0.5 min}{%
  Four steps. Replicate Smirnov and Thurner's own checks, to confirm the data
  behaves the way they reported. Run centrality on its own, with no
  popularity control --- the naive test. Add popularity, sweep phi, and pool
  all five groups --- the test we planned last time. And check the
  assumptions behind that test --- where we found one had failed. Here's what
  actually happened at each step.
  \Cue{Four steps, fast; land on "here's what happened."}
}

\scriptframe{The replication holds}{about 0.8 min}{%
  First check: does a student's own past GPA predict their next GPA? Yes,
  very strongly --- this is by far the largest coefficient in the whole
  regression. Does a friend's past GPA add anything on top of that? No ---
  not once you control for the student's own trajectory. Second check: are
  new friendships more GPA-similar than friendships about to end? In the
  university data, yes, clearly. In school, the same direction, but not
  significant, probably because the school sample is smaller. Both of these
  match Smirnov and Thurner's original findings, now in this new panel
  setting with a real time lag. That's the foundation everything else in
  this talk builds on.
  \Cue{Own-GPA strong, friend-GPA null; new-vs-dropped; "matches prior work" as foundation.}
}

\scriptframe{The naive effect is popularity, relabelled}{about 0.8 min}{%
  Now the first genuinely new test. Katz-Bonacich centrality on its own, net
  of friend GPA, pooled across all five groups: positive, and highly
  significant. Looks like position matters. But add raw in-degree --- just
  how many people like you --- to that same regression, and the centrality
  coefficient drops to essentially zero, no longer significant. Meanwhile
  in-degree itself is large and highly significant. That's the tell: centrality
  and popularity are highly correlated in these networks, so whichever one
  you leave out of the regression looks like it's doing all the work. So the
  naive centrality effect wasn't really about network position at all. It
  was popularity, relabelled.
  \Cue{Naive positive; add degree, it vanishes; degree itself is what was doing the work.}
}

\scriptframe{The headline test}{about 0.5 min}{%
  This is the test we actually planned last time. Popularity and friend GPA
  are both controlled. Phi is swept across six values, pooled across all
  five groups. At phi equals 0.85 --- the value we fixed in advance as our
  headline --- we get a small, negative, statistically significant
  coefficient. Read at face value, that says network position has a small,
  independent, negative effect on grades, net of popularity and friend
  quality.
  \Cue{The sweep table briefly; the fixed headline number; "read at face value" signals more is coming.}
}

\scriptframe{But we found a gap in it}{about 0.6 min}{%
  Here's the problem we found. That regression was supposed to control for a
  student's own past GPA. But remember: GPA is only measured repeatedly for
  the school group. For all four university cohorts, it's a single static
  number in the data, not something that varies over time. You literally
  cannot build an own-past-GPA regressor for over eighty percent of that
  pooled sample. The model we promised last time had that control written
  into it. The test we actually ran did not have it.
  \Cue{Restate the promised control; the static-GPA fact; the honest admission, in the brief's own words.}
}

\scriptframe{The corrected test}{about 1.4 min}{%
  So we restricted the confirmatory test to the one valid sample. That's the
  school group, with a real time lag and a real own-past-GPA control. It
  comes to two thousand and eighty-eight observations, from five hundred
  thirty-five students. Before trusting it, we checked
  whether that own-GPA control collides with the other regressors. It
  doesn't. The correlations are all small, and the variance inflation factor
  for own-GPA is 1.11 --- textbook territory. Katz and in-degree still
  correlate highly with each other, around 0.87, but that's expected, and
  it's exactly what sweeping phi across the whole range is designed to
  handle. The result: at every single value of phi, the coefficient is not
  significant, and the sign isn't even consistent from one end of the dial
  to the other. Put the three tests side by side and you get three different
  pictures. The naive test, positive. The pooled test, negative. The
  corrected test, looking like nothing at all. Net of popularity, friend
  GPA, and a student's own past performance, once the test is actually
  specified correctly, we find no evidence that network position
  independently predicts grades.
  \Cue{Sample and validity; the VIF check briefly; all six phi values null; the three-way picture; the headline.}
}

\begin{samepage}
\scriptframe{Conclusion}{about 0.9 min}{%
  So here's the answer. Net of popularity, friend GPA, and a student's own
  past performance, there's no evidence that network position predicts
  grades on its own. Why does that still matter? Because the progression itself is the finding --- a
  naive test that looked positive, a pooled test that looked negative, and a
  correctly specified test that looks like nothing, all from the same data.
  Pooling groups with mismatched data structures can manufacture a
  significant result that the one correctly specified subsample simply
  doesn't support. Practically, that means interventions built around
  widening students' network reach aren't supported by what we've found
  here. Next, we're testing whether that null result comes from weak,
  one-way ``likes,'' by looking only at mutually confirmed ties.
  \Cue{Headline answer; the progression as the real finding; implication; one-line forward pointer. No "thank you."}
}
\end{samepage}

\ScriptAppendix

\scriptframe{Sample and network construction}{if asked}{%
  Give the Harvard Dataverse source and CC0 licence, the roughly 6,200
  Russian students, the node counts by cohort, and the directed ``like'' tie
  definition with its 3-month windows. If asked why ``likes'' rather than
  self-reported friendship: name the 24--27\% reciprocity rate directly, and
  connect it to the reciprocal-ties test planned next.
}

\scriptframe{Robustness: the full \(\phi\) range}{if asked}{%
  Walk through the two lines on the chart. Before the own-GPA fix, pooled
  swings from significantly positive at low phi to significantly negative at
  high phi. After the fix, school-only stays flat and non-significant across
  the same range. The point isn't that one phi value is right --- it's that
  the corrected line never leaves the null region anywhere on the dial.
}

\scriptframe{Checking for multicollinearity}{if asked}{%
  If someone worries that adding a student's own past GPA should collide
  with the other right-hand-side variables: walk through the correlation
  table --- 0.008 with Katz, 0.065 with in-degree, 0.122 with out-degree,
  0.285 with friend GPA --- and the VIF of 1.11. That correlation with the
  outcome is the whole point of a control variable; what would be a problem
  is correlation with the \emph{other regressors}, and that's what this table
  rules out.
}

\scriptframe{What's next: reciprocal ties}{if asked}{%
  Only 24 to 27 percent of ``like'' ties go both ways. The next step
  recomputes centrality using only those mutually confirmed ties. That tests
  directly whether a one-way ``like'' is too weak a tie for the model's
  mechanism to operate through, and whether the null result changes once
  ties are restricted to ones both people confirm.
}

\scriptframe{Related literature}{if asked}{%
  Calv{\'o}-Armengol, Patacchini and Zenou (2009) is the model itself, tested
  on a single cross-section. Smirnov and Thurner (2017) is the paper and data
  we extend. Giulietti, Vlassopoulos and Zenou (2020) is a different peer
  effects design, using reference groups rather than friendship ties, by the
  same course instructor. Positioning: this is the first time centrality has
  been tested with a real time lag and an own-performance control, which the
  original cross-sectional data could not support.
}

\Needspace{7\baselineskip}
\section*{Standing Q\&A stance}
Hear the whole question and count to three. Answer its strongest version in
one sentence before adding detail. If asked why we're presenting a null
result: say plainly that the progression --- naive, pooled, corrected --- is
itself the finding, and that the corrected test is the only honest one.

\Needspace{8\baselineskip}
\section*{Live cue layer}
\raggedright\small
\begin{enumerate}\setlength{\itemsep}{0.45em}
  \item \textbf{\KeyQuestion} --- two antecedents, the gap each leaves, our combined plan.
  \item \textbf{Recap: design and data} --- the equation, each control's job, the static-GPA fact.
  \item \textbf{What we did} --- four steps, fast.
  \item \textbf{The replication holds} --- own-GPA strong, friend-GPA null, new-vs-dropped.
  \item \textbf{The naive effect is popularity, relabelled} --- positive, then vanishes with degree.
  \item \textbf{The headline test} --- the sweep, the fixed 0.85 headline, "at face value."
  \item \textbf{But we found a gap in it} --- the promised control that wasn't actually there.
  \item \textbf{The corrected test} --- valid sample, VIF check, all-phi null, the three-way picture.
  \item \textbf{Conclusion} --- the answer, the progression as the finding, the forward pointer.
\end{enumerate}

\end{document}
